\question[4]
Analysiere den folgenden Python-Code und gib an, welche Ausgabe er erzeugt.
\begin{lstlisting}
a = [2, 4, 9, 12, 1, 5]
b = []
teilsumme = 0
for x in a:
    teilsumme += x
    b.append(teilsumme)
print(b)
\end{lstlisting}

\begin{solutionbox}{5cm}

Bei jedem Durchlauf wird das aktuelle Element x zu teilsumme
addiert; der aktuelle Wert von teilsumme wird danach an b
angehängt (laufende/kumulierte Summe).

\begin{lstlisting}
x=2:  teilsumme=2  -> b=[2]
x=4:  teilsumme=6  -> b=[2, 6]
x=9:  teilsumme=15 -> b=[2, 6, 15]
x=12: teilsumme=27 -> b=[2, 6, 15, 27]
x=1:  teilsumme=28 -> b=[2, 6, 15, 27, 28]
x=5:  teilsumme=33 -> b=[2, 6, 15, 27, 28, 33]

Ausgabe: [2, 6, 15, 27, 28, 33]
\end{lstlisting}
\end{solutionbox}